\chapter{Introduction}

\section{Overview and Motivation}
Acoustic monitoring provides a non-invasive, cost-effective method for predictive maintenance in industrial environments. By analyzing the sound emitted by rotating machinery (pumps, fans, and motors), it is possible to detect early-stage faults like bearing wear or imbalance before catastrophic failure occurs. This project develops an automated system to distinguish between "Normal" and "Abnormal" acoustic signatures across multiple machine types.

\section{System Objectives}
To deliver a production-ready solution, the system must address several technical constraints:
\begin{itemize}
    \item \textbf{Noise Robustness:} Handle varying industrial background noise levels.
    \item \textbf{Data Foundations:} Standardizing and augmenting a dataset of 56,000+ recordings.
    \item \textbf{Leakage Prevention:} Using clustering-aware validation to ensure generalization.
    \item \textbf{Class Imbalance:} Handling datasets where "Normal" operation is heavily dominant.
    \item \textbf{Model Benchmarking:} Optimizing both XGBoost and EfficientNet architectures.
\end{itemize}

\section{Scope and Context}
The system addresses a 6-class classification problem covering the operational states of three distinct machines. The resulting pipeline serves as a blueprint for a real-time monitoring service capable of providing maintenance teams with instant anomaly alerts.
